\documentclass[11pt]{article}

\usepackage{amsmath}
\usepackage{amssymb}
\usepackage{graphicx}
\usepackage{hyperref}
\usepackage{geometry}
\geometry{margin=1in}

\title{GRA-Swarm-Engine: Resonance-Based Swarm Intelligence for Emergent Multi-Agent Optimization}

\author{
Oleg Bitsoev\\ 
Independent Researcher \\
\texttt{https://github.com/qqewq}
}

\date{2026}

\begin{document}

\maketitle

\begin{abstract}

This work introduces \textbf{GRA-Swarm-Engine}, an experimental multi-agent framework designed to investigate emergent collective intelligence through resonance-driven swarm dynamics.

The system combines concepts from swarm intelligence, evolutionary algorithms, and resonance-based state amplification. Agents explore a shared solution space while interacting through a General Resonance Algorithm (GRA) that promotes coherent agent states.

We hypothesize that large-scale resonant swarms may produce emergent problem-solving behaviors and serve as a foundation for decentralized artificial intelligence architectures.

\end{abstract}

\section{Introduction}

Traditional artificial intelligence systems rely on large centralized models trained on massive datasets. While effective, this paradigm differs significantly from biological intelligence, which often emerges from distributed interacting components.

Swarm intelligence provides an alternative computational paradigm in which simple agents interact locally to produce global intelligent behavior.

In this work we introduce \textbf{GRA-Swarm-Engine}, a research framework that explores how:

\begin{itemize}
\item multiple agents explore a shared state space
\item agent states interact through resonance
\item swarm dynamics amplify coherent solutions
\item collective intelligence emerges from decentralized interactions
\end{itemize}

The project provides a minimal experimental implementation of this architecture.

\section{System Architecture}

The system consists of a swarm of agents:

\[
s_i \in S
\]

where:

\begin{itemize}
\item \(s_i\) represents the state of agent \(i\)
\item \(S\) is the solution or hypothesis space
\end{itemize}

Each iteration of the swarm performs the following steps:

\begin{enumerate}
\item agents generate candidate states
\item resonance between agents is evaluated
\item agents receive scores based on resonance
\item the swarm evolves via selection and mutation
\end{enumerate}

This process produces a dynamic adaptive swarm capable of exploring complex solution spaces.

\section{Resonance Mechanism}

The core idea of the GRA-Swarm system is the resonance between agent states.

Given two agents \(i\) and \(j\):

\[
R(i,j) = similarity(s_i, s_j)
\]

In the current implementation, similarity is measured using cosine similarity:

\[
R(i,j) =
\frac{s_i \cdot s_j}
{\|s_i\| \|s_j\|}
\]

The total resonance score for an agent is:

\[
A_i = \sum_{j} R(i,j)
\]

Agents with higher resonance scores are considered more aligned with the swarm.

\section{Swarm Evolution}

The swarm evolves according to a simple evolutionary mechanism.

First, agents are ranked by resonance score:

\[
A_1, A_2, ..., A_n
\]

The top-performing agents are selected:

\[
S_{top} = \{s_1, s_2, ..., s_k\}
\]

New agents are generated through mutation:

\[
s'_i = s_i + \epsilon
\]

where \(\epsilon\) is a small stochastic perturbation.

This evolutionary process allows the swarm to continuously explore the solution space.

\section{Collective Dynamics}

As the swarm evolves, several collective phenomena may appear:

\begin{itemize}
\item clustering of agent states
\item convergence toward stable solutions
\item distributed exploration of the search space
\item adaptive swarm behavior
\end{itemize}

These behaviors resemble properties observed in biological swarms and decentralized systems.

\section{Collective Memory}

The architecture may include a collective memory component that stores high-quality states discovered by the swarm.

Such memory can accumulate:

\begin{itemize}
\item best-performing states
\item discovered patterns
\item effective strategies
\end{itemize}

Collective memory allows the swarm to retain useful information across iterations.

\section{LLM Swarm Extension}

The GRA-Swarm architecture can be extended with language model agents.

In this configuration:

\begin{enumerate}
\item agents generate ideas using language models
\item ideas are evaluated using resonance metrics
\item the swarm selects and refines promising ideas
\end{enumerate}

This approach enables a form of distributed reasoning where multiple agents collaboratively explore a solution space.

\section{Experiments}

The GRA-Swarm-Engine repository contains several experimental modules:

\begin{itemize}
\item swarm optimization simulation
\item resonance clustering experiments
\item emergent role experiments
\item prototype LLM swarm reasoning
\end{itemize}

These experiments demonstrate basic swarm dynamics and provide a foundation for further research.

\section{Future Work}

Several research directions remain open:

\begin{itemize}
\item large-scale swarms with thousands of agents
\item adaptive resonance functions
\item meta-optimization of swarm parameters
\item integration with advanced language models
\item decentralized AGI architectures
\end{itemize}

Exploring these directions may reveal new approaches to distributed intelligence.

\section{Conclusion}

GRA-Swarm-Engine provides a minimal research framework for investigating resonance-based swarm intelligence.

Instead of relying on a single centralized model, the system explores how collective intelligence can emerge from many interacting agents.

While preliminary, this architecture suggests that decentralized resonance-based systems may provide a promising direction for future AI research.

\section*{Code}

Implementation:

\[
\texttt{https://github.com/qqewq/GRA-Swarm-Engine}
\]

\end{document}